% !PRINTED 292
\weberpage{292}

Nach (3) k\"onnen wir die Compositionsreihen der beiden Gruppen $P/Q$ und $P/R$
daraus herleiten, die wir so andeuten wollen:
\weq{5}{P/Q,\ Q'/Q,\ Q''/Q,\ \ldots,\ 1}
\weq{6}{P/R,\ Q'/R,\ Q''/R,\ \ldots,\ Q/R \ldots}
Nun ist der Index des Theilers $Q'/Q$ von $P/Q$ gleich dem Index des Theilers $Q'$
von $P$, also auch gleich dem Index des Theilers $Q'/R$ von $P/R$ (\S.\,6, 1.), und
Gleiches gilt von den folgenden Gliedern. Wir sprechen also den Satz aus:

\noindent 3. \gesperrt{Sind $Q$ und $R$ Normaltheiler von $P$, und ist $R$ ein
Theiler von $Q$, so ist die Indexreihe von $P/Q$ ein Theil der Indexreihe von
$P/R$..} % [sic] double period as printed

Sind die Indices $j_1, j_2, \ldots, j_\nu$ lauter Primzahlen, so ist die L\"osung
der Resolvente $\chi(y) = 0$ auf die L\"osung einer Kette von cyklischen
Gleichungen der Grade $j_1, j_2, \ldots, j_\nu$ reducirbar; diese Resolvente ist
also metacyklisch in dem Sinne, wie wir diesen Begriff im siebzehnten Abschnitte
des ersten Bandes festgestellt haben, wonach die metacyklischen Gleichungen mit den
sonst algebraisch l\"osbar genannten identisch sind.

Eine irreducible Gleichung $f(x) = 0$ ist metacyklisch, wenn die Indexreihe ihrer
Gruppe aus lauter Primzahlen besteht. Wir haben an der erw\"ahnten Stelle die
Bedingungen f\"ur metacyklische Gleichungen von Primzahlgrad untersucht, und
m\"ussen jetzt diese Betrachtungen f\"ur den allgemeinen Fall durchf\"uhren, dass
der Grad der Gleichung beliebig zusammengesetzt ist.

\begin{center}
\S.\,74.

\medskip

\gesperrt{Metacyklische Gleichungen.}
\end{center}

\bigskip

Wenn $f(x) = 0$ eine irreducible Gleichung $m^{\mathrm{ten}}$ Grades und $P$ ihre
\gesperrt{Galois}'sche Gruppe ist, so ist $P$ transitiv. Eine Compositions- und Indexreihe
f\"ur $P$ sei
\weq{1}{
\begin{gathered}
P,\ P_1,\ P_2,\ \ldots,\ P_{\mu-1},\ 1 \\
j_1,\ j_2,\ \ldots,\ j_{\mu-1},\ j_\mu.
\end{gathered}
}
Es wird, da die letzte Gruppe $1$ intransitiv ist, in der Compositionsreihe einmal
eine intransitive Gruppe auftreten, und es sei also $P_\lambda$ die erste
intransitive Gruppe der Reihe (1). Dann
%CONT>>
